\chapter{Methylmalonic Acid}
\section*{What Is This Test?}
Methylmalonic acid (MMA) is a dicarboxylic acid intermediate in the metabolic pathway of propionyl-CoA, derived from the catabolism of branched-chain amino acids (isoleucine, valine, methionine, threonine), odd-chain fatty acids, and cholesterol. The conversion of methylmalonyl-CoA (as D-methylmalonyl-CoA, converted to L-methylmalonyl-CoA by methylmalonyl-CoA racemase) to succinyl-CoA is catalyzed by methylmalonyl-CoA mutase (MUT), an enzyme that requires adenosylcobalamin (adenosyl-B12, the mitochondrial coenzyme form of vitamin B12) as an essential cofactor. In vitamin B12 (cobalamin) deficiency, MUT activity is impaired, leading to the accumulation of methylmalonyl-CoA, which is hydrolyzed to methylmalonic acid and excreted in the urine and blood. Serum MMA is the single most sensitive and specific metabolic marker for vitamin B12 deficiency at the tissue level \cite{savage1994sensitivity}. Serum MMA is elevated in vitamin B12 deficiency BEFORE the development of anemia or macrocytosis, making it the gold standard for detecting early or subclinical B12 deficiency \cite{stabler2013vitamin}. MMA is normal in folate deficiency (because the MUT enzyme requires vitamin B12, not folate), making it useful for differentiating B12 deficiency from folate deficiency when both homocysteine and B12 levels are borderline. MMA is elevated in: vitamin B12 deficiency (cobalamin), pernicious anemia, gastric bypass, ileal resection, and the rare inborn error of metabolism methylmalonic acidemia (MUT deficiency, severe neonatal metabolic acidosis) \cite{carmel2000current}.
\section*{Alternative Names}
MMA; Methylmalonic Acid; Serum Methylmalonic Acid; Urine Methylmalonic Acid
\section*{Principle}
Serum MMA is measured by LC-MS/MS (liquid chromatography-tandem mass spectrometry) or gas chromatography-mass spectrometry (GC-MS) \cite{tietz2012textbook}. Serum specimen should be collected in an SST (gold-top) tube. MMA levels are increased in renal insufficiency because MMA is renally cleared; moderate elevation (0.4--1.5 mcmol/L) may be secondary to CKD rather than B12 deficiency.
\section*{History}
In the 1950s--1960s, studies by Flavin and colleagues showed that adenosylcobalamin is the coenzyme required for the conversion of methylmalonyl-CoA to succinyl-CoA. In 1967, methylmalonic acidemia was first described by Oberholzer and colleagues in infants with severe neonatal acidosis and developmental delay. In 1968, Rosenberg and colleagues demonstrated that some forms of methylmalonic acidemia respond to vitamin B12 therapy, distinguishing B12-responsive from B12-unresponsive disease. The development of gas chromatography--mass spectrometry in the 1970s permitted sensitive measurement of serum and urine MMA, and in the 1980s--1990s Lindenbaum and colleagues established elevated MMA and homocysteine as the most sensitive markers of tissue-level cobalamin deficiency \cite{lindenbaum1990diagnosis}. Today, LC-MS/MS is the reference method for serum MMA measurement.
\section*{Why This Test Is Ordered}
\begin{itemize}
  \item Confirmation of suspected vitamin B12 deficiency when serum B12 is borderline (200--300 pg/mL) or low-normal, especially in patients with normal MCV or unexplained neurologic symptoms
  \item Differentiation of vitamin B12 deficiency from folate deficiency (MMA is elevated only in B12 deficiency; homocysteine is elevated in both)
  \item Evaluation of macrocytic anemia, peripheral neuropathy, or subacute combined degeneration of the spinal cord when B12 status is uncertain
  \item Monitoring of B12 status after gastric bypass, ileal resection, or during pernicious anemia therapy
  \item Workup of newborn or infant metabolic acidosis with hyperammonemia to evaluate methylmalonic acidemia (MUT deficiency)
  \item Assessment of B12 status in pregnancy, older adults, and strict vegans, where early detection before anemia is desired
\end{itemize}
\section*{Primary vs Secondary Test}
Methylmalonic acid is a secondary, confirmatory test. The primary tests for B12 deficiency are serum vitamin B12 and the CBC with MCV; MMA (with homocysteine) is ordered when these are borderline or equivocal. For methylmalonic acidemia, plasma and urine MMA are primary diagnostic tests in the metabolic workup.
\section*{How This Test Is Performed}
A blood sample is collected by venipuncture into a serum separator (gold-top) tube, and a urine sample may be collected in a plain container. Serum or urine MMA is measured by LC-MS/MS or gas chromatography--mass spectrometry. Because MMA is renally cleared, serum creatinine is measured simultaneously to interpret the result. Results are typically available within 1--3 days.
\section*{What Preparation Is Needed}
None required; MMA can be measured in a fasting or non-fasting sample. However, the sample should be collected before starting B12 replacement, because treatment rapidly normalizes MMA and obscures the diagnosis. Serum creatinine should be checked at the same time for interpretation.
\section*{Contraindications and Precautions}
No contraindications to venipuncture. Interpretive cautions: (1) renal insufficiency elevates MMA (0.4--1.5+ mcmol/L) independent of B12 status; (2) MMA falls within days of B12 therapy, so testing during or after treatment underestimates deficiency; (3) MMA may be mildly elevated in small bowel bacterial overgrowth, hypovolemia, and, in newborns, by physiologic immaturity; (4) MMA is not useful for monitoring B12 replacement once stores are replete.
\section*{Risks}
Minimal risk. Slight pain or bruising at the venipuncture site. Rarely: hematoma, infection, or vasovagal syncope.
\section*{What Happens After the Test}
An elevated MMA with a low-normal B12 confirms tissue-level cobalamin deficiency and warrants B12 replacement, with MMA rechecked after 1--3 months. If MMA is elevated but B12 is normal and renal function is preserved, intrinsic factor antibody testing is performed to evaluate pernicious anemia. In infants, markedly elevated MMA with metabolic acidosis prompts urgent metabolic evaluation and treatment for methylmalonic acidemia.
\section*{Understanding Results}
\subsection*{Normal (<0.30--0.40 mcmol/L)}
B12 status adequate at the tissue level. MUT enzyme activity is normal. Normal methylmalonic acid.
\subsection*{Elevated (>0.40 mcmol/L)}
Vitamin B12 deficiency at the tissue level. MMA >0.40 mcmol/L has >95\% sensitivity and specificity for B12 deficiency after adjusting for renal function. Elevated MMA in a patient with normal serum B12 (200--300 pg/mL) indicates tissue-level B12 deficiency and is a strong indication for B12 replacement therapy. Renal insufficiency: MMA may be elevated (0.4--1.5+ mcmol/L) due to impaired renal clearance alone, without B12 deficiency. Severe elevation (>1.0--2.0+ mcmol/L with normal renal function) indicates significant B12 deficiency. Methylmalonic acidemia: extremely elevated MMA (>100--1,000+ mcmol/L) in neonates with metabolic acidosis, hyperammonemia, and developmental delay.
\section*{Normal Results}
Serum MMA: <0.30--0.40 mcmol/L (varies by laboratory, upper limit typically 0.30--0.37 mcmol/L). Mildly elevated in CKD (0.4--1.5 mcmol/L). B12 deficiency: >0.40 mcmol/L. MMA is the best single test for confirming B12 deficiency.
\section*{Follow-up Tests}
Serum vitamin B12, serum homocysteine (also elevated in B12 deficiency), serum folate, intrinsic factor antibody and anti-parietal cell antibody (pernicious anemia), complete blood count (macrocytosis, anemia), and Schilling test (historical, no longer performed in most centers).
\section*{References}
\bibliographystyle{unsrtnat}
\bibliography{references}

\clearpage
\section*{Regulatory Status and Authorization Date}
Regulatory status applies to the specific assay, instrument, manufacturer, and intended use rather than to the test name alone. No single approval date can be assigned to this test category. For a commercial assay, verify the current product in the relevant regulator's database: the U.S. Food and Drug Administration (FDA) clearance, approval, or authorization; India's Central Drugs Standard Control Organisation (CDSCO) licence or permission; the European Union In Vitro Diagnostic Regulation (EU IVDR) CE certificate issued through the applicable conformity-assessment route; or the United Kingdom Medicines and Healthcare products Regulatory Agency (MHRA) registration and UKCA/CE status. Laboratory-developed tests may not have an individual FDA, CDSCO, EU IVDR, or MHRA approval date.
\clearpage
